\section{The same original vertical generic stalk from every component point}

Let $Y$ be integral. For an actual specialization $\eta\leadsto p$ in $Y$,
write $S=\mathcal O_{Y,p}$, $D=\mathcal O_{Y,\eta}$ and let
$j:S\longrightarrow D$ be the original specialization map of the structure
sheaf. No localization, ring isomorphism, valuation comparison, or equality
of native points is an input.

\begin{definition}[The original generization coefficient map]
\label{def:gauss-common-original-generization-map}
\lean{EnvelopeProperty.GaussCommonVerticalStalks.generizationMap}
\leanok
The coefficient map is exactly $j=\mathcal O_Y(\eta\leadsto p)$.
Put $P=j^{-1}(\mathfrak m_D)$, with its derived prime-ideal structure.
\end{definition}
\begin{proof}
Use the actual structure-sheaf specialization homomorphism, with its
induced algebra structure on $D$. The inverse image of the genuine maximal
ideal is prime. This construction does not choose an affine presentation
or a substitute local map.
\end{proof}

\begin{lemma}[The actual generization stalk is the actual localization]
\label{lem:gauss-common-original-generization-localization}
\uses{def:gauss-common-original-generization-map}
\lean{EnvelopeProperty.GaussCommonVerticalStalks.generization_isLocalization}
\leanok
With the original coefficient action $j$, $D$ is the localization of the
SAME $S$ at $P$.
\end{lemma}
\begin{proof}
Choose one actual affine neighborhood $V$ of $p$. Specialization places
$\eta$ in the same $V$. Put $A=\Gamma(V,\mathcal O_Y)$ and retain the
actual germ maps $g_p:A\to S$ and $g_\eta:A\to D$.
The structure-sheaf specialization identity gives
$j g_p=g_\eta$. The existing affine-stalk localization theorem makes
$D$ a localization of $A$ at the prime corresponding to $\eta$.

Every element outside $P$ has image outside $\mathfrak m_D$, hence has a
unit image in $D$. Given $d\in D$, its actual affine-stalk fraction
representation gives $a,b\in A$ with $g_\eta(b)$ a unit and
$d g_\eta(b)=g_\eta(a)$. Thus $g_p(b)\notin P$ and
$d j(g_p(b))=j(g_p(a))$, which is the required fraction representation
over $S$. Finally, both original stalks inject into the SAME integral
function field. Composition of specialization maps shows that the map
from $S$ to that field factors through $j$, so $j$ is injective. Equality
of two images therefore already implies equality of the two original
coefficients; the localization equality axiom holds with denominator $1$.
These three actual properties prove the asserted localization structure.
\end{proof}

Now let $i:C\hookrightarrow Y$ be a genuine closed immersion with $C$
integral, let $x\in C$, $p=i(x)$ and $\eta=i(\eta_C)$. The original
specialization $\eta_C\leadsto x$ gives the original
$\eta\leadsto p$. Suppose the actual local equation $\pi\in S$ is prime
and satisfies $\ker(i_x^\#)=(\pi)$. This is a geometric equation for the
chosen component, not a supplied valuation, center, or ring comparison.

\begin{lemma}[The true component equation identifies the actual generization prime]
\label{lem:gauss-common-original-component-prime}
\uses{def:gauss-common-original-generization-map}
\lean{EnvelopeProperty.GaussCommonVerticalStalks.component_generizationPrime}
\leanok
For this SAME original generization map, $P=(\pi)$.
\end{lemma}
\begin{proof}
The actual closed-immersion stalk map
\[
 i_{\eta_C}^\#:D\to\mathcal O_{C,\eta_C}=C.\mathrm{functionField}.
\]
It is surjective. Its kernel is therefore the genuine maximal ideal of $D$.
Naturality of the actual stalk maps for $\eta_C\leadsto x$ gives
\[
 i_{\eta_C}^\# j
 =\bigl(\mathcal O_{C,x}\longrightarrow C.\mathrm{functionField}\bigr)i_x^\#.
\]
The displayed stalk-to-function-field map is injective because $C$ is
integral. Consequently the kernel of the displayed composite is exactly
$\ker(i_x^\#)$. The left-hand description is precisely
$j^{-1}(\mathfrak m_D)=P$, and the original component equation identifies
the right-hand kernel with $(\pi)$.
\end{proof}

\begin{theorem}[The same vertical DVR is the same original global generic stalk]
\label{thm:gauss-common-original-component-stalk}
\uses{lem:gauss-common-original-generization-localization,lem:gauss-common-original-component-prime,def:gauss-surface-original-residual-vertical-rings,lem:gauss-surface-original-residual-vertical-dvrs}
\lean{EnvelopeProperty.GaussCommonVerticalStalks.componentStalkEquiv,EnvelopeProperty.GaussCommonVerticalStalks.componentStalkEquiv_coefficient,EnvelopeProperty.GaussCommonVerticalStalks.componentStalk_isDVR}
\leanok
There is an internally constructed actual ring isomorphism
\[
 E_x:S_{(\pi)}\xrightarrow{\sim}\mathcal O_{Y,i(\eta_C)}.
\]
For EVERY original $s\in S$ it satisfies
$E_x(s/1)=j(s)$. If this SAME $S$ is Noetherian, then the SAME original
generic stalk $\mathcal O_{Y,i(\eta_C)}$ is an actual DVR.
Neither an isomorphism nor a DVR structure on that generic stalk is
supplied as an input.
\end{theorem}
\begin{proof}
The preceding two lemmas make the original generic stalk a localization
of $S$ at the prime complement of the SAME $(\pi)$. The old original
vertical ring is the literal localization at that same multiplicative
set. The localization universal property gives the canonical
$S$-algebra isomorphism between these two actual rings. Its algebra
coefficient law is exactly the stated formula, retaining every original
coefficient and the actual specialization map.

Under the genuine Noetherian hypothesis, the existing original vertical
ring theorem derives that $S_{(\pi)}$ is a DVR from the original nonzero
prime parameter. Transfer that structure through the constructed whole
ring isomorphism. No normalization theorem or additional regularity,
codimension, valuation, or native-point certificate is needed here.
\end{proof}

This checkpoint identifies the actual rings and maps needed for a common
vertical point. It does not yet assert equality of normalized native
points, equality of their whole function-field norms, a common active
family for source convexity, or a residual degree inequality.
